% Originally: content/week-01.tex, \section{From Analysis I to several variables}
% Source: Maarten Cnoops/Woche 1.pdf, p. 1
\section{From Analysis I to several variables}
\label{sec:analysis1_recap}

Several of the definitions you will meet in the coming chapters are exactly the definitions you
already know from \emph{Analysis I}, only stripped of their reliance on the real line. A
\newterm{sequence} (\germanterm{Folge}) $(x_k)_{k \in \mathbb{N}}$ in a set $X$ converging to
$x \in X$, a function being continuous at a point, a set being open or closed --- all of these
notions carry over to $\mathbb{R}^n$ and, as \cref{sec:structured_spaces} will show, to much more
general spaces than $\mathbb{R}^n$. What does \emph{not} carry over automatically is anything
that relied on the \emph{order} of $\mathbb{R}$: there is no \qt{$\leq$} on $\mathbb{R}^n$ for
$n \geq 2$, so statements like the intermediate value theorem or \qt{every bounded sequence has a
convergent subsequence} need new proofs, or new hypotheses, once we leave $\mathbb{R}$.

\begin{remark}
That last fact --- bounded does \emph{not} imply convergent, but bounded sequences in
$\mathbb{R}^n$ do always have a convergent \emph{subsequence} --- is exactly the statement of the
Bolzano--Weierstrass theorem, revisited via \cref{ex:1.1} below. It
reappears as the Heine--Borel theorem once compactness is introduced in \cref{sec:compactness}.
\end{remark}

% Supplement: Linus Lüchinger/Slides/Slides-02-23.pdf
\begin{ainote}
Linus Lüchinger's first exercise-session slides collect the most common mistakes on the first
problem sheet, whose problems are distributed through this chapter. Two are worth flagging before
attempting \cref{ex:1.1}: on part 2, most students first guessed $f(x) = \log|x| + c$ for a single
constant $c$ --- wrong, since the two constants on $(-\infty,0)$ and $(0,\infty)$ may genuinely
differ (see \cref{sec:prerequisites_solutions}). On part 4, most students assumed a function with
everywhere-positive second derivative cannot have a global maximum; it can, at the boundary of the
domain.
\end{ainote}

% Source: exercises/Ex1_Analysis2_eng.pdf, p. 1
\begin{exercise}[1.1 --- Recap Questions]
\label{ex:1.1}
\begin{enumerate}[label=\textbf{\arabic*.}]
  \item Let $\{x_k\}_{k \in \mathbb{N}} \subset [1,2] \subset \mathbb{R}$ be a sequence. Is it
    necessarily convergent?
  \item Describe all differentiable functions $f : \mathbb{R} \setminus \{0\} \to \mathbb{R}$ such
    that $f'(x) = 1/x$ for all $x \neq 0$.
  \item Let $f : \mathbb{R} \to \mathbb{R}$ be continuous and bounded. Is it true that
    \[
      \int_0^{\pi/2} f(\cos(x))\,dx = \int_0^1 \frac{f(t)}{\sqrt{1-t^2}}\,dt \, ?
    \]
  \item Let $f : [0,\infty) \to \mathbb{R}$ have $f''(x) > 0$ for all $x \geq 0$. Can $f$ have a
    global maximum point?
  \item Let $\phi : \mathbb{R}^2 \to \mathbb{R}^2$ be a linear map and let
    $Q := [0,1] \times [0,1]$ be the unit square. How does $\phi(Q)$ look like? What are the
    possibilities?
  \item Consider the vector space $V$ of polynomials of degree less than $4$, and the function
    $D : V \to V$ which associates to a polynomial $p$ its derivative $p'$. Is $D$ a linear map?
    If so, what is the dimension of its null space?
\end{enumerate}
\end{exercise}
